% !TeX root = ../main.tex
\section{Решение задачи №1 ex1}
\vspace{1em}
\textbf{1. Исходные данные}
\begin{itemize}
    \item Мощность шума на выходе БВЧ: $P_{\text{вых}} = 20 \cdot 10^{-9} \text{ Вт}$
    \item Шумовая полоса: $\Pi_{\text{ш}} = 2 \text{ МГц} = 2 \cdot 10^6 \text{ Гц}$
    \item Коэффициент усиления проходной мощности: $K_{p\text{(дБ)}} = 60 \text{ дБ}$
    \item Температура эквивалента антенны (стандартная): $T_0 = 290 \text{ К}$
    \item Постоянная Больцмана: $k = 1.38 \cdot 10^{-23} \text{ Вт/(Гц}\cdot\text{К)}$
\end{itemize}

\vspace{1em}
\textbf{2. Перевод коэффициента усиления мощности}

Переведем заданный в децибелах коэффициент усиления мощности БВЧ в относительные единицы (разы):
\begin{align*}
K_p &= 10^{\frac{K_{p\text{(дБ)}}}{10}} = 10^{\frac{60}{10}} = 10^6 \text{ раз}
\end{align*}

\vspace{1em}
\textbf{3. Расчет суммарной мощности шума на входе}

Зная мощность на выходе и коэффициент усиления, найдем общую суммарную мощность шума, приведенную ко входу приемника:
\begin{align*}
P_{\text{вх.}\Sigma} &= \frac{P_{\text{вых}}}{K_p} = \frac{20 \cdot 10^{-9}}{10^6} = 20 \cdot 10^{-15} \text{ Вт}
\end{align*}

\vspace{1em}
\textbf{4. Расчет мощности шума эквивалента антенны}

Используем формулу Найквиста для определения мощности теплового шума, поступающего от согласованной антенны:
\begin{align*}
P_A &= k \cdot T_0 \cdot \Pi_{\text{ш}} \\
P_A &= 1.38 \cdot 10^{-23} \cdot 290 \cdot 2 \cdot 10^6 \\
P_A &= 800.4 \cdot 10^{-17} = 8.004 \cdot 10^{-15} \text{ Вт}
\end{align*}

\vspace{1em}
\textbf{5. Определение мощности собственного шума приёмника}

Мощность суммарного приведенного ко входу шума складывается из мощности шума антенны и мощности собственного приведенного шума приёмника:
\begin{align*}
P_{\text{вх.}\Sigma} &= P_A + P_{\text{вх.собств}}
\end{align*}

Отсюда выразим искомую мощность собственного шума приемника:
\begin{align*}
P_{\text{вх.собств}} &= P_{\text{вх.}\Sigma} - P_A \\
P_{\text{вх.собств}} &= 20 \cdot 10^{-15} - 8.004 \cdot 10^{-15} = 11.996 \cdot 10^{-15} \text{ Вт}
\end{align*}

\textbf{Ответ:} Мощность приведенного ко входу собственного шума приемника составляет $11.996 \cdot 10^{-15} \text{ Вт}$ (или приблизительно $1.2 \cdot 10^{-14} \text{ Вт}$).